\scanpage{602}

\mspage{11}

que $\unsure{\mathcal U|U}_1$ contienne \unsure{comme|\deleted{comme} un} élément un ens. infini.

Si $\mathcal U=\mathcal U_1$, et si \deleted{\unsure{on a $k=\mathbb Q$ est}} \inserted{$k$ est un corps} infini, p. ex. $k=\unsure{\mathbb Q|Q}$, alors pour \unsure{tt} \inserted{\unsure{ou}} $\forall E\in\mathcal U$ et pour $\unsure{\mathrm{tt}|\forall}$ ultrafiltre \uline{non} \uline{trivial} $\Phi$ sur $E$, $k_E$ est une extension \uline{\deleted{\unsure{tri}}non triviale} de $k$. Mais si $\mathcal U$\inserted{$=\mathcal U_\alpha$} est un univers quelconque, \deleted{je présume} \inserted{et} $\unsure{\pi|\kappa}$ le premier cardinal \deleted{\unsure{tel q}} \inserted{$\in\mathcal U$} qui n'appartienne pas \unsure{en|à} un univers $\mathcal U'\subsetneq\mathcal U$ strictement plus petit, je présume que pour $\forall E\in\mathcal U$, \deleted{\unsure{tout}} $\inserted{\pi\text{-}}$ultrafiltre sur $E$ est trivial.

\uline{Application} \uline{à un exemple} \uline{d'une} \uline{sous-catégorie} \uline{non accessible}.

Soit $\unsure{\mathcal M|M}$ une catégorie $(\mathcal U\text{-})$accessible, et $a\in\operatorname{Ob}\mathcal M$, d'où un foncteur \deleted{\unsure{covariant}}

\[
\begin{aligned}
f^o\colon\mathcal M^o&\unsure{\longmapsto|\longrightarrow}\mathrm{Ens}\\
x&\longmapsto\operatorname{Hom}_{\mathcal M}(x,a)\overset{\mathrm{\unsure{déf|def}}}{=}a(x),
\end{aligned}
\]

foncteur commutant aux \deleted{\unsure{petites}} $\underset{\longleftarrow}{\lim}$ quelconques, \deleted{\unsure{donc}} tel que $f\colon\mathcal M\longrightarrow(\mathrm{Ens})^o$ commute aux $\underset{\longrightarrow}{\lim}$ quelconques, et en particulier $f$ est accessible. Considérons la sous-catégorie \inserted{str.} pleine

% archive sheet number: 601
